\section{The many-row seed}\label{sec:manyrow}

\begin{theorem}\label{thm:manyrow}
There are absolute constants $B,c_*,C_*>0$ such that the following holds. Let
$2\le d$, $n\ge Bd^2$, let $X\in\R^{n\times d}$ have $\norm{x_i}^2=a=d/n$ for all
$i$, and suppose $\eta:=\opn{X^TX-I}$ satisfies $0<\eta\le c_*$. Then there is
an ENP frame $W$ with $\fro{X-W}^2\le C_*\eta d$.
\end{theorem}

Throughout, $S=X^TX$, so
$\opn{S-I}\le\eta\le c_*\le\frac12$ (the choice below gives $c_*\le\frac1{64}$), and
$e_i=x_i/\sqrt a$.

\subsection{The noise}

In the Frobenius space $\R^{n\times d}$ consider the subspaces
\[
 \mathcal E=\{Z:\ip{x_i}{z_i}=0\ \forall i\},\qquad
 \mathcal F=\{Z\in\mathcal E:\ X^TZ+Z^TX=0\},
\]
and let $m=\dim\mathcal E-\dim\mathcal F\le\dim\mathrm{Sym}_d\le d^2\le n$. For a
standard Gaussian $g\in\R^{n\times d}$ put
\[
 Z_0=n^{-1/2}\Pi_{\mathcal E}g,\qquad Z=n^{-1/2}\Pi_{\mathcal F}g,\qquad
 Z_\perp=Z_0-Z=n^{-1/2}\Pi_{\mathcal E\ominus\mathcal F}g .
\]
$Z$ and $Z_\perp$ are independent. The rows $z_{0,i}=n^{-1/2}(I-e_ie_i^T)g_i$
of $Z_0$ are independent. Since $\Cov(Z)\preceq\Cov(Z_0)$, each row $z_i$ of
$Z$ is Gaussian with covariance $\Sigma_i\preceq n^{-1}(I-e_ie_i^T)$, and
$\Cov(\mathrm{vec}\,Z)\preceq I/n$. Moreover
\begin{equation}\label{eq:mr-noise}
 \E\fro Z^2=\frac{n(d-1)-m}{n}\le d,\qquad
 \E\fro{Z-Z_0}^2=\frac mn\le\frac{d^2}n .
\end{equation}

\begin{lemma}[Row moments]\label{lem:rowmoments}
Let $r_i=\norm{z_i}^2/a$, $\mu_i=\E r_i$ and $\delta_i=1-\mu_i$. Then
$1/d\le\delta_i\le1$, $\sum_ia\delta_i=1+m/n\le2$, and
\[
 \mathcal M:=a\sum_i\E\bigl[(r_i-1)^2(1+r_i)^2\bigr]\le C_{\mathcal M},
\]
with $C_{\mathcal M}$ absolute.
\end{lemma}

\begin{proof}
$\tilde\Sigma_i=\Sigma_i/a\preceq d^{-1}I$ has trace $\mu_i\le(d-1)/d$, so
$\tr\tilde\Sigma_i^2\le1/d$ and $\tr\tilde\Sigma_i^4\le d^{-3}$. For
$\xi_i=r_i-\mu_i$ the Gaussian moment formulas give $\E\xi_i^2=2\tr\tilde\Sigma_i^2\le
2/d$ and $\E\xi_i^4=12(\tr\tilde\Sigma_i^2)^2+48\tr\tilde\Sigma_i^4\le60/d^2$.
Next $\sum_ia\delta_i=d-\E\fro Z^2=1+m/n$ by~\eqref{eq:mr-noise}. Finally
$(r_i-1)^2\le2\xi_i^2+2\delta_i^2$ and $(1+r_i)^2\le8+2\xi_i^2$, so
\[
 \E(r_i-1)^2(1+r_i)^2\le16\E\xi_i^2+4\E\xi_i^4+16\delta_i^2+4\delta_i^2\E\xi_i^2
 \le\frac{C}d+C\delta_i^2 ,
\]
and $a\sum_i(C/d+C\delta_i^2)\le C+Ca\sum_i\delta_i\le3C$.
\end{proof}

\subsection{Row normalisation and the exact second-order identity}

For $0<t\le1$ set
\[
 v_i=\frac{x_i+tz_i}{\sqrt{1+t^2r_i}},\qquad
 v_{0,i}=\frac{x_i+tz_{0,i}}{\sqrt{1+t^2r_{0,i}}},\qquad
 r_{0,i}=\norm{z_{0,i}}^2/a .
\]
Since $z_i\perp x_i$, $\norm{v_i}^2=\norm{v_{0,i}}^2=a$ exactly. The map
$y\mapsto\sqrt a\,y/\norm y$ is the metric projection onto the closed ball of
radius $\sqrt a$ on $\{\norm y\ge\sqrt a\}$, hence $1$-Lipschitz there. Thus
\begin{equation}\label{eq:mr-contract}
 \fro{V-X}\le t\fro Z,\quad \fro{V-V_0}\le t\fro{Z-Z_0},\quad
 \opn V\le\opn{X+tZ},\quad\opn{V_0}\le\opn{X+tZ_0},
\end{equation}
the last two because $V=\Diag(\lambda)(X+tZ)$ with $0<\lambda_i\le1$, and similarly
for $V_0$.

Expanding $1/(1+t^2r_i)=1-t^2c_i$ with
\[
 c_i=\frac{r_i}{1+t^2r_i},
\]
and using $X^TZ+Z^TX=0$, one finds the exact identity
\begin{equation}\label{eq:mr-identity}
 V^TV=S+t^2Q(Z)+R_3+R_4,\qquad Q(Z)=Z^TZ-\sum_ir_ix_ix_i^T,
\end{equation}
\[
 R_3=-t^3\sum_ic_i(x_iz_i^T+z_ix_i^T),\qquad
 R_4=-t^4\sum_ic_i(z_iz_i^T-r_ix_ix_i^T).
\]

\begin{lemma}[Mean and fluctuation of $Q$]\label{lem:Q}
$\opn{\E Q(Z)-(I-S)}\le m/n$ and $\E\fro{Q(Z)-\E Q(Z)}^2\le8d^2/n$.
\end{lemma}

\begin{proof}
For $Z_0$: $\E Z_0^TZ_0=n^{-1}\sum_i(I-e_ie_i^T)=I-S/d$ and $\E r_{0,i}=(d-1)/d$,
so $\E Q(Z_0)=I-S/d-(1-1/d)S=I-S$. As $Z_0=Z+Z_\perp$ with independent centred
summands and $Q$ is quadratic, $\E Q(Z)=\E Q(Z_0)-\E Q(Z_\perp)$, and
$\E Q(Z_\perp)$ is the difference of the positive semidefinite matrices
$\E Z_\perp^TZ_\perp$ and $\sum_i\E(\norm{z_{\perp,i}}^2/a)x_ix_i^T$, both of trace
$\E\fro{Z_\perp}^2=m/n$; so its norm is at most $m/n$.

Entry $(\alpha,\beta)$ of $Q$ is the quadratic form $\sum_iz_i^TB_{\alpha\beta,i}z_i$
with $B_{\alpha\beta,i}=\frac12(E_{\alpha\beta}+E_{\beta\alpha})-(e_i)_\alpha(e_i)_\beta I$.
For a Gaussian vector with covariance $\Sigma\preceq I/n$ and symmetric
coefficient matrix $\mathsf B$, the variance of the quadratic form is
$2\fro{\Sigma^{1/2}\mathsf B\Sigma^{1/2}}^2\le2\fro{\mathsf B}^2/n^2$. Since
$\sum_{\alpha,\beta}\fro{B_{\alpha\beta,i}}^2\le2d^2+2d\le4d^2$ for each $i$,
summing over $\alpha,\beta$ and the $n$ blocks gives the claim.
\end{proof}

\begin{lemma}[Remainder]\label{lem:remainder}
Let $\bar c=1/(1+t^2)$, $\Gamma=\Diag(c_i-\bar c)$ and
$\Lambda=\Diag(c_ir_i-\bar c)$. Then
\[
 \opn{R_3}\le2t^3\opn X\fro{\Gamma Z},\qquad
 \opn{R_4}\le t^4\bigl(\opn Z^2+\opn Z\fro{\Gamma Z}+\opn S+\opn X\fro{\Lambda X}\bigr),
\]
and $\E\fro{\Gamma Z}^2\le\mathcal M$, $\E\fro{\Lambda X}^2\le\mathcal M$.
\end{lemma}

\begin{proof}
Because $\sum_i(x_iz_i^T+z_ix_i^T)=X^TZ+Z^TX=0$, the constant part $\bar c$ of
the weights drops out of $R_3$:
$R_3=-t^3(X^T\Gamma Z+Z^T\Gamma X)$, which gives the first bound. Next
$\sum_ic_iz_iz_i^T=\bar cZ^TZ+Z^T\Gamma Z$ and $\sum_ic_ir_ix_ix_i^T=\bar cS+X^T\Lambda X$,
with $\opn{Z^T\Gamma Z}\le\opn Z\fro{\Gamma Z}$ and similarly for $\Lambda$, and
$\bar c\le1$. For the moments,
\[
 c_i-\bar c=\frac{r_i-1}{(1+t^2r_i)(1+t^2)},\qquad
 c_ir_i-\bar c=\frac{(r_i-1)(1+r_i+t^2r_i)}{(1+t^2r_i)(1+t^2)},
\]
so $|c_i-\bar c|\le|r_i-1|$ and, as $1+r+t^2r\le(1+r)(1+t^2r)$,
$|c_ir_i-\bar c|\le|r_i-1|(1+r_i)$. Hence
$\fro{\Gamma Z}^2=a\sum_i(c_i-\bar c)^2r_i\le a\sum_i(r_i-1)^2(1+r_i)^2$ and
$\fro{\Lambda X}^2=a\sum_i(c_ir_i-\bar c)^2\le a\sum_i(r_i-1)^2(1+r_i)^2$; take
expectations and use \cref{lem:rowmoments}.
\end{proof}

\begin{remark}
The termwise bound $\opn{R_3}\le2t^3\sum_ic_i\norm{x_i}\norm{z_i}\approx2t^3d$
would only allow $\eta\lesssim d^{-2}$. Centring the weights at
$\bar c$, which is allowed by the first-order constraint, replaces the $n$
terms of size $a$ by fluctuations of relative size $d^{-1/2}$, whose
$\ell_2$-sum over the rows is $O(1)$.
\end{remark}

\subsection{A dense reference graph}

\begin{lemma}\label{lem:mr-dense}
There are absolute $h,t_0\in(0,1]$ and $c_1>0$ such that for $d\ge2$ and
$t\le t_0$, with
probability at least $1-ne^{-c_1(n-1)}$, every $i$ has at least $0.9(n-1)$
indices $j\ne i$ with
\[
 |\ip{v_{0,i}}{v_{0,j}}|\ge h\,t\sqrt{a/n}.
\]
\end{lemma}

\begin{proof}
Write $z_{0,j}/\sqrt a=d^{-1/2}g'_j$ with $g'_j\sim N(0,I-e_je_j^T)$
independent. Fix $i$ and condition on row $i$; let $f=v_{0,i}/\sqrt a$, a unit
vector. For $j\ne i$,
\[
 \frac{\ip{v_{0,i}}{v_{0,j}}}{a}=\frac{\nu_j}{N_j},\qquad
 \nu_j=u_j+\tfrac t{\sqrt d}\ip f{g'_j},\quad u_j=\ip f{e_j},\quad
 N_j=\bigl(1+\tfrac{t^2}d\norm{g'_j}^2\bigr)^{1/2}\ge1,
\]
and the target inequality reads $|\nu_j|/N_j\ge ht/\sqrt d$. By Markov,
$\Prob(N_j>2)\le t^2/3$. On $\{N_j\le2\}$ failure requires
$|\nu_j|<2ht/\sqrt d$. Here $\nu_j\sim N(u_j,t^2(1-u_j^2)/d)$. If $|u_j|\le\frac12$
the variance is at least $3t^2/(4d)$ and the Gaussian density bound gives
probability at most $4h/\sqrt{3\pi/2}\le3h$. If $|u_j|>\frac12$ and $2ht\le\frac14$,
failure requires $\frac t{\sqrt d}|\ip f{g'_j}|>\frac14$, of probability at most
$16t^2/d$. Choose $h$ and $t_0$ so that $3h+t_0^2/3+8t_0^2\le\frac1{50}$ (this also gives
$2ht\le\frac14$, and $16t^2/d\le8t^2$ as $d\ge2$). The
failure indicators for $j\ne i$ are conditionally independent with
probabilities at most $\frac1{50}$, so by the exponential Markov inequality
(with parameter $1$) more than $0.1(n-1)$ failures have probability at most
$\exp(-(n-1)[0.1-\frac{e-1}{50}])\le e^{-c_1(n-1)}$. Take a union bound over
$i$.
\end{proof}

\subsection{Proof of \texorpdfstring{\cref{thm:manyrow}}{Theorem}}

\emph{Constants.} Fix $h,t_0,c_1$ from \cref{lem:mr-dense}. Put $\kappa_0=h^2/200$,
$C'=36^2\cdot20$, and
$C_{\rm mix}=125/h^2+3$, then $\Theta=\max\{42,C_{\rm mix}\}$ and
$\theta=\min\{1/(4\Theta),h/60\}$. Let $C_R=300+100\sqrt{C_{\mathcal M}}$, choose $t_1=\min\{t_0,\theta/(16C_R)\}$, and set
\[
 t^2=4\eta/\theta,\qquad c_*=\theta t_1^2/4,
\]
so that $\eta\le c_*$ gives $t\le t_1$. Finally choose $B$ so large that
$n\ge Bd^2$ implies $d^2/n\le\theta/4$, $\sqrt{160d^2/n}\le\theta/4$,
$2aC'd/\kappa_0\le\frac14$ (this bounds $\sum_{i\in B_0}2a$ for the set $B_0$
below, which has $|B_0|\le C'd/\kappa_0$), $n\ge30$, and $ne^{-c_1(n-1)}\le0.01$.

\emph{The event.} By Markov's inequality, \eqref{eq:mr-noise},
\cref{lem:Q,lem:remainder}, the net bound (\cref{lem:gauss}(c)) for $Z$ and
$Z_0$, and \cref{lem:mr-dense}, with positive probability all of the following
hold:
\[
 \fro Z^2\le20d,\quad \fro{Z-Z_0}^2\le\frac{20d^2}n,\quad
 \fro{Q-\E Q}^2\le\frac{160d^2}n,\quad
 \fro{\Gamma Z}^2\le20\mathcal M,\quad \fro{\Lambda X}^2\le20\mathcal M,
\]
$\opn Z,\opn{Z_0}\le16$, and the conclusion of \cref{lem:mr-dense}: the five Markov
events fail with probability at most $\frac1{20}$ each, the two net bounds
(\cref{lem:gauss}(c) with $\sigma^2=1/n$, $u=16$) with probability at most
$2\cdot5^{n+d}e^{-8n}$ each, and \cref{lem:mr-dense} with probability at most
$0.01$. Fix such a sample.

\emph{Spectral error.} By~\eqref{eq:mr-identity},
$V^TV-I=(1-t^2)(S-I)+t^2(\E Q-(I-S))+t^2(Q-\E Q)+R_3+R_4$. Since
$\opn X\le\sqrt2$, $\opn S\le\frac32$ and $t\le1$, \cref{lem:remainder} and the
event give
$\opn{R_3}+\opn{R_4}\le t^3\bigl(2\sqrt{40\mathcal M}+256+16\sqrt{20\mathcal M}
+\tfrac32+\sqrt{40\mathcal M}\bigr)\le C_Rt^3$, and therefore
\[
 \delta:=\opn{V^TV-I}\le\eta+t^2\frac{d^2}n+t^2\sqrt{\frac{160d^2}n}+C_Rt^3
 \le\frac\theta4t^2+\frac\theta4t^2+\frac\theta4t^2+\frac\theta{16}t^2<\theta t^2 .
\]

\emph{Graph.} By~\eqref{eq:mr-contract}, $\opn V,\opn{V_0}\le\sqrt2+16\le18$
and
$\fro{VV^T-V_0V_0^T}\le36\,t\fro{Z-Z_0}$, so
$\fro{VV^T-V_0V_0^T}^2\le C't^2d^2/n$. Let $B_0$ be the set of rows whose
contribution to this sum exceeds $\kappa_0at^2$; then $|B_0|\le C'd/\kappa_0$. A row
$i\notin B_0$ has at most $\kappa_0at^2/((h/2)^2t^2a/n)=n/50$ indices with
$|(VV^T-V_0V_0^T)_{ij}|\ge\frac h2t\sqrt{a/n}$, so by \cref{lem:mr-dense} it has at
least $0.85n$ indices $j\ne i$ with $|(VV^T)_{ij}|\ge\frac h2t\sqrt{a/n}$.

Let $U=V(V^TV)^{-1/2}$, $P=UU^T$. By \cref{lem:align} (rows of $P-VV^T$) a row loses at most
$6a\theta^2t^4/((h/4)^2t^2a/n)\le n/20$ of these indices, so each $i\notin
B_0$ has at least $0.8n$ indices $j\ne i$ with $P_{ij}^2\ge(h^2/16)at^2/n$.
Also $a/2\le P_{ii}\le2a$ and $\sum_{i\in B_0}P_{ii}\le\frac14$.
\Cref{lem:core} with $\mathcal B=B_0$, part (i) with $\vartheta=0.3$ and
$\gamma=h^2at^2/16$, and part (iii) with $\alpha=a$ (so $R=\frac8{3a}$,
$F_{\mathcal B}\le\frac43$) shows that
$H(L_P)\le\frac{7/3}{0.3\gamma}+\frac8{3a}\le C_{\rm mix}/(at^2)$.

\emph{Seed.} $\fro{U-X}^2\le2d\delta^2+2t^2\fro Z^2\le42t^2d$, the diagonal error
is at most $2a\delta\le2\theta at^2\le at^2/(2\Theta)$, and the barrier constant
is at most $\Theta/(at^2)$. So $U$ is a $(\Theta,t)$-seed, and
\cref{cor:seed} gives an ENP frame within $3\Theta t^2d=(12\Theta/\theta)\eta d$
of $X$.
\qed
